\bookchapter{الفصل 67: الصقل والنشر — الاختبار والتنقيح والنشر على App Store}

\chapterintro{أنجزت التطبيق، وماذا بعد؟ يجيب هذا الفصل عن ثلاث مسائل في الطريق «من كود يعمل» إلى «منتج قابل للتسليم»: كيف تنقح بفاعلية لتجد الأخطاء، وكيف تكتب اختبارات تمنع الانتكاس، وكيف ترفع التطبيق إلى App Store ليحمّله العالم. والنشر ليس مجرد طقس — بل هو الفاصل الذي يحول عملاً إلى منتج.}

\booksection{أولًا: أدوات التنقيح الثلاث — print ونقاط التوقف والمعاينة}

\begin{center}
\begin{tabularx}{\linewidth}{|X|X|X|}
\hline
\textbf{الأداة} & \textbf{وظيفتها} & \textbf{كيف تستخدمها} \tabularnewline
\hline
print + وحدة التحكم & عرض سجلات التشغيل & في الكود اكتب print("القيمة=\textbackslash{}(x)")، وانظر الناتج في وحدة التحكم أسفل Xcode \tabularnewline
التنقيح بنقاط التوقف & إيقاف البرنامج عند سطر معين ومعرفة قيم المتغيرات & انقر بجوار رقم السطر فيظهر علام زرقاء، فيتوقف التشغيل عنده، وتتقدم خطوة خطوة \tabularnewline
معاينة SwiftUI + المنقح & مشاهدة الواجهة وتسلسل العروض لحظيًا & لوحة \#Preview + فاحص View Hierarchy أثناء التشغيل \tabularnewline
\hline
\end{tabularx}
\end{center}

\begin{center}\small\color{gray}أدوات التنقيح الثلاث\end{center}

\begin{tcbinfo}{نصيحة: كيف تقرأ تقرير الانهيار؟}
عند انهيار Swift تطبع وحدة التحكم “Fatal error” مع السبب، ويتوقف Xcode عند السطر المخطئ مبرزًا إياه. أما الانهيارات التاريخية على الجهاز الحقيقي فيمكن رؤيتها في Xcode Window Organizer Crashes (يتطلب حسابًا مدفوعًا). واجعل أولوية الكلمات المفتاحية في البحث “الجملة الأولى من سبب الخطأ + swiftui” — فبالسؤال الدقيق تأتي الإجابة الدقيقة.

\end{tcbinfo}

\booksection{ثانيًا: الاختبار — تأمين على الكود}

الاختبار اليدوي لا يُعتمد — فتصلح ميزة فتتعطل أخرى بصمت. الاختبار «دواء الندم» و«السترة المضادة للرصاص» معًا. يستخدم iOS إطار XCTest وينقسم إلى نوعين:

\begin{center}
\begin{tabularx}{\linewidth}{|X|X|X|X|}
\hline
\textbf{المستوى} & \textbf{ماذا يختبر} & \textbf{Target} & \textbf{السرعة} \tabularnewline
\hline
اختبارات الوحدات & المنطق الصافي (الحساب والتحقق) & MyFirstAppTests & بالثواني \tabularnewline
اختبارات الواجهة & سلوك الواجهة الحقيقي & MyFirstAppUITests & بالدقائق \tabularnewline
الاختبار اليدوي & إحساس وتجربة إنسان حقيقي & —— & لا غنى عنه أبدًا \tabularnewline
\hline
\end{tabularx}
\end{center}

\begin{center}\small\color{gray}مستويات الاختبار في iOS (المشروع النموذجي يتضمن مسبقًا هدفَي اختبار Test Target)\end{center}

\begin{lstlisting}[language=Swift]
import XCTest

// استخرج منطق "تحويل نص المبلغ إلى رقم" كدالة نقية
func parseAmount(_ input: String) -> Double? {
    guard let v = Double(input.trimmingCharacters(in: .whitespaces)),
          v > 0 else { return nil }
    return v
}

final class AmountTests: XCTestCase {
    func testEmptyInputReturnsNil() {
        XCTAssertNil(parseAmount(""))
    }
    func testNegativeReturnsNil() {
        XCTAssertNil(parseAmount("-5"))
    }
    func testNormalValueParses() {
        XCTAssertEqual(parseAmount("12.5"), 12.5)
    }
}
\end{lstlisting}

تشغيل الاختبارات: CmdU. وبعد أن يصير كل شيء أخضر تعدل الكود وقلبك مطمئن. والتحقق من المبالغ وحسابات الملخص في دفتر الحسابات ينبغي أن تحميها اختبارات وحدات.

\booksection{ثالثًا: التوقيع والتغليف — ختم على التطبيق}

كل تطبيق iOS يجب أن يوقع حتى يمكن تثبيته — إثباتًا أن «هذه الحزمة منك ولم يمسها تغيير». العلاقة الهرمية:

\begin{itemize}
  \item مرحلة التعلم: Apple ID مجاني + التوقيع التلقائي في Xcode (Personal Team)، يعمل على الجهاز الحقيقي ويجدد كل 7 أيام؛
  \item مرحلة التوزيع: حساب مطور مدفوع (99 دولارًا سنويًا) يصدر الشهادات وملفات الوصف، وعندها فقط يمكن النشر على المتجر والتثبيت طويل الأمد.
\end{itemize}

التغليف: من قائمة Xcode اختر Product Archive (الأرشفة)، وبعد انتهاء البناء افتح Organizer واختر Distribute App ثم App Store Connect ثم ارفع. يجب أن يكون هدف الأرشفة Generic iOS Device (لا المحاكي)، وهي خطوة ينساها المبتدئون كثيرًا.

\textbf{تنبيه لتجنب الفخاخ: ضياع الشهادة = انقطاع التحديثات}\par

بمجرد ضياع شهادة التوقيع والمفتاح الخاص لن تستطيع إصدار تحديثات للتطبيق القديم — وسيضطر المستخدمون إلى تثبيت تطبيق «جديد كليًا». احتفظ بنسخة احتياطية من الشهادة (بتصدير .p12 من حلقة المفاتيح) واحفظ كلمة السر جيدًا. ولحسن الحظ توفر Apple «التوقيع التلقائي» وإدارة الشهادات السحابية (فعّل Automatically manage signing في Xcode)، وهو ما يجنبك معظم فخاخ العمل اليدوي.

\booksection{رابعًا: الدورة الكاملة للنشر على App Store}

تسجيل حساب مطوّر

تجهيز مواد المتجر

رفع حزمة البناء

ملء معلومات الخصوصية والمراجعة

مراجعة Apple (1-7 أيام)

النشر

\begin{center}\small\color{gray}شكل 5: ست خطوات للنشر على App Store\end{center}

\begin{itemize}
  \item التسجيل: على developer.apple.com، حساب فردي بـ 99 دولارًا سنويًا؛
  \item المواد: الأيقونة (1024×1024 دون زوايا دائرية — فApple تقصها بنفسها)، واللقطات (بمقاسات iPhone المختلفة)، والوصف، ورابط سياسة الخصوصية؛
  \item استبيان الخصوصية: أعلن بصدق ما جمعته من بيانات — دفتر الحسابات غير المتصل بالشبكة يمكنه اختيار «لا نجمع أي بيانات»؛
  \item المراجعة: أول إرسال ينتظر عادة من يوم إلى سبعة، وإذا رُفض فعدّل بندًا بندًا وفق «ملاحظات المراجعة» — فأسباب الرفض من Apple تكتب عادة بتفصيل محدد، فنفذها كما هي؛
  \item TestFlight: قبل النشر أرسل النسخة التجريبية إلى أصدقائك للاختبار الداخلي (حتى 10000 شخص)، واجمع جولة من الملاحظات قبل الإرسال الرسمي، فتبدو احترافية تمامًا.
\end{itemize}

\booksection{خامسًا: كيف تصنع مواد المتجر — الانطباع الأول هو عدد التنزيلات}

صفحة المتجر هي واجهة المحل التي يراها المستخدم أثناء «تجوله». ثلاث مواد تحسم الانطباع الأول:

\begin{itemize}
  \item الأيقونة: 1024×1024، ولا تكتفي بلقطة شاشة حلاً عجلاً. إن لم تملك مهارات التصميم فولّد شكلًا بسيطًا بأدوات الذكاء الاصطناعي، والمعيار أن تميزها من نظرة واحدة حتى بعد تصغيرها على سطح المكتب؛
  \item اللقطات: مجموعتان على الأقل (بمقاسات iPhone السائدة) تعرضان أهم مشاهد الوظائف؛ وكل لقطة معها جملة بسيطة مفهومة («سجل عملية في ثانية» أفضل بعشرة أضعاف من «حل متقدم لإدارة الحسابات»)؛
  \item الوصف: الجملة الأولى توضح «ما هذا + لمن هو»، ولا تكدد صفات مثل «الأكثر احترافية والأقوى» — فالصدق أنجع من المبالغة.
\end{itemize}

\begin{tcbpractice}{هيا نجرّب: «سلّم» دفتر الحسابات مرة واحدة}
\booksection{سادسًا: قائمة مختصرة للإبلاغ عن الانهيارات وتحسين الأداء}

قبل نشر التطبيق، يستحق هذان الأمران نصف ساعة لكل منهما:

\booksubsection{مراقبة الانهيارات}

عدم انهياره في اختباراتك لا يعني أنه لن ينهار على هواتف المستخدمين. توفر Apple تقارير انهيار مجانية في Xcode Organizer (وحساب مدفوع لرؤية سجلات انهيار الأجهزة الحقيقية)؛ والحل الأريح أكثر هو دمج SDK من طرف ثالث (مثل Firebase Crashlytics). والنشر دون مراقبة للانهيارات كالقيادة والعينان مغمضتان.

\booksubsection{أربعة بنود لفحص الأداء}

\begin{enumerate}
  \item القوائم بـ LazyVStack أو List: لا ترص القائمة الطويلة بـ VStack العادي — فالتحميل الكسول هو الدرس الأول في الأداء؛
  \item لا تدس الصور الأصلية كما هي: استخدم AsyncImage (للصور البعيدة) والتخزين المؤقت للنظام، فصورة أصلية بحجم 5MB تكفي لتختنق بها الأجهزة القديمة؛
  \item لا تنجز أعمالًا ثقيلة على الخيط الرئيسي: الشبكة والملفات الكبيرة كلها عبر async/await؛
  \item قس باستخدام Instruments: محلل الأداء المدمج في Xcode (Time Profiler وفحص تسريب الذاكرة) — قس أولًا ثم حسن، ولا تعتمد على الحدس.
\end{enumerate}

(1) أضف ثلاث اختبارات وحدات لمنطق parseAmount ونظائره (وشغّلها بـ CmdU حتى تخضرّ)؛ (2) اصطاد بنفسك بنقطة توقف انهيارًا صنعته عمدًا وأصلحه؛ (3) أنشئ حزمة عبر Product Archive وادعُ أصدقاءك لاختبارها عبر TestFlight؛ (4) إن وسعت سجل حساب مطور وقدّم دفتر الحسابات للنشر فعليًا. الخطوات الثلاث الأولى يمكن إنجازها الليلة.

\end{tcbpractice}

\begin{tcbsummary}{خلاصة الفصل}
\begin{itemize}
  \item أدوات التنقيح الثلاث: سجلات print، ونقاط التوقف، ومعاينة SwiftUI وفاحص العروض
  \item XCTest: اختبارات وحدات (CmdU) + اختبارات واجهة + اختبار يدوي، والمشروع النموذجي يتضمنها مسبقًا
  \item التوقيع: حساب مجاني بتوقيع تلقائي (يجدد كل 7 أيام)؛ والتوزيع لا يتم إلا بحساب مدفوع
  \item التغليف: Generic iOS Device + Archive + Distribute؛ ولا بد من نسخة احتياطية للشهادة
  \item النشر في ست خطوات: التسجيل المواد الرفع استبيان الخصوصية المراجعة النشر؛ وTestFlight للاختبار الداخلي أولًا
\end{itemize}

\end{tcbsummary}
